\section{Problemas resueltos de regresión lineal simple}

\subsection*{Enunciados de los problemas}

% Recordar
\begin{problema}
	\label{prob:ab5bf4a}
	Se estudia la relación entre las horas de estudio ($X$) y la calificación obtenida ($Y$) en un examen de estadística modular. Se recolectaron datos aleatorios de $n = 10$ estudiantes:
	\begin{center}
		\begin{tabular}{|l|c|c|c|c|c|c|c|c|c|c|}\hline
			\textbf{Horas ($X$)} & 2 & 3 & 4 & 5 & 6 & 7 & 8 & 9 & 10 & 12 \\\hline
			\textbf{Calificación ($Y$)} & 45 & 50 & 55 & 60 & 65 & 70 & 75 & 80 & 85 & 90 \\\hline
		\end{tabular}
	\end{center}
	\begin{enumerate}
		\item Calcule los estimadores de mínimos cuadrados ordinarios (MCO) para el intercepto ($\hat{\beta}_0$) y la pendiente ($\hat{\beta}_1$).
		\item Escriba la ecuación empírica del modelo de regresión lineal ajustado $\hat{Y} = \hat{\beta}_0 + \hat{\beta}_1 X$.
		\item Prediga la calificación esperada para un estudiante que dedica 7.5 horas al estudio antes del examen.
		\item Calcule el coeficiente de determinación $R^2$ e interprete su significado en el contexto pedagógico.
	\end{enumerate}
\end{problema}

% Comprender
\begin{problema}
	\label{prob:e5cf043}
	Se analiza la relación estructural entre el peso de paquetes postales ($X$, en kg) y el costo final de envío ($Y$, en dólares) en una empresa de logística. Un software estadístico reporta el siguiente modelo empírico basado en una muestra de $n = 20$ envíos con una masa media de $\bar{X} = 10.0$ kg:
	\begin{align}
		\hat{Y} = 5.0 + 2.5X, \quad R^2 = 0.85, \quad \text{SCT} = 200.0.
	\end{align}
	\begin{enumerate}
		\item Interprete las estimaciones numéricas del intercepto ($\hat{\beta}_0 = 5.0$) y la pendiente ($\hat{\beta}_1 = 2.5$).
		\item Calcule la predicción puntual para el costo de envío de un paquete cuyo peso es exactamente $15.0$ kg.
		\item A partir de la Suma de Cuadrados Total ($\text{SCT} = 200.0$) y el coeficiente $R^2$, determine la Suma de Cuadrados de la Regresión ($\text{SCR}$) y la Suma de Cuadrados del Error Residual ($\text{SCE}$).
		\item Calcule la varianza muestral del error de estimación $\hat{\sigma}^2$ y el error estándar residual del modelo ($S_e$).
	\end{enumerate}
\end{problema}

% Aplicar
\begin{problema}
	\label{prob:5a642fc}
	Un ingeniero de redes eléctricas urbanas analiza si las fluctuaciones de la temperatura ambiente exterior ($X$, en °C) impactan de manera estadísticamente probada el consumo total diario de energía en la red de distribución ($Y$, en megavatios-hora, MWh). Para un período de $n = 15$ días estivales, se obtuvieron las métricas descriptivas:
	\begin{align}
		\bar{X} = 25.0 \text{ °C}, \quad \bar{Y} = 450.0 \text{ MWh}, \quad S_{XX} = 500.0, \quad S_{XY} = 2000.0.
	\end{align}
	\begin{enumerate}
		\item Calcule los parámetros estimados por mínimos cuadrados $\hat{\beta}_0$ y $\hat{\beta}_1$ y exprese la ecuación de regresión.
		\item Si la Suma de Cuadrados Total en el consumo eléctrico fue $\text{SCT} = 10000.0$, calcule la Suma de Cuadrados Explicada por la Regresión ($\text{SCR}$) y la Suma de Cuadrados del Error Residual ($\text{SCE}$).
		\item Realice la prueba formal de significancia sobre la pendiente ($H_0: \beta_1 = 0$ vs. $H_a: \beta_1 \neq 0$) a un nivel de significancia $\alpha = 0.05$, calculando el estadístico de prueba $t = \dfrac{\hat{\beta}_1}{SE(\hat{\beta}_1)}$ donde $SE(\hat{\beta}_1) = \sqrt{\dfrac{\text{CME}}{S_{XX}}}$, y verifique su congruencia con el cuantil crítico $t_{0.025, \, 13} = 2.160$.
	\end{enumerate}
\end{problema}

% Analizar
\begin{problema}
	\label{prob:e97453b}
	\textbf{Deducción analítica de las ecuaciones normales de Gauss e insesgo exacto en MCO:} considere el modelo estadístico lineal de regresión poblacional $Y_i = \beta_0 + \beta_1 X_i + \varepsilon_i$, $i=1,\dots,n$, donde $X_i$ es no aleatoria con $S_{XX} > 0$, y los errores satisfacen $E[\varepsilon_i] = 0$, $\text{Var}(\varepsilon_i) = \sigma^2$ y $\text{Cov}(\varepsilon_i, \varepsilon_j) = 0$ para $i \neq j$.
	\begin{enumerate}
		\item Definiendo $Q(\beta_0, \beta_1) = \sum_{i=1}^n (Y_i - \beta_0 - \beta_1 X_i)^2$, aplique las condiciones de primer orden ($\partial Q/\partial\beta_0 = 0$, $\partial Q/\partial\beta_1=0$) para deducir el sistema de \textbf{ecuaciones normales de Gauss}.
		\item A partir de la resolución del sistema normal, obtenga las fórmulas cerradas $\hat{\beta}_1 = \dfrac{\sum(X_i-\bar{X})(Y_i-\bar{Y})}{\sum(X_i-\bar{X})^2}$ y $\hat{\beta}_0 = \bar{Y}-\hat{\beta}_1\bar{X}$.
		\item Aplicando linealidad de la esperanza y sustituyendo el modelo verdadero de $Y_i$, demuestre que $\hat{\beta}_1$ y $\hat{\beta}_0$ son estimadores estrictamente \textbf{insesgados}.
	\end{enumerate}
\end{problema}

% Evaluar
\begin{problema}
	\label{prob:5c55133}
	Un ingeniero metalúrgico estudia la resistencia a la fatiga cíclica en megaciclos ($Y$) de una aleación estructural en función del contenido porcentual de titanio adicionado ($X$). El modelo calibrado con $n = 18$ especímenes arrojó $\hat{Y} = 120.0 + 15.0X$, $\bar{X} = 3.0\%$, $S_{XX} = 40.0$, $S_e = 4.2$ megaciclos.
	\begin{enumerate}
		\item Calcule el error estándar del valor poblacional esperado $\text{SE}(\hat{\mu}_{Y|X_0}) = S_e\sqrt{1/n + (X_0-\bar{X})^2/S_{XX}}$ y el error estándar de predicción individual $\text{SE}(Y_{\text{pred}|X_0}) = S_e\sqrt{1+1/n+(X_0-\bar{X})^2/S_{XX}}$, ambos en $X_0=6.0\%$ ($t_{0.025,16}=2.120$).
		\item Construya el intervalo de confianza del 95\% para la respuesta media y el intervalo de predicción del 95\% para un espécimen individual en $X_0=6.0\%$. Evalúe y explique por qué el intervalo de predicción resulta sustancialmente más ancho que el de confianza de la media, incluso cuando $n\to\infty$.
	\end{enumerate}
\end{problema}

% Crear
\begin{problema}
	\label{prob:58d3723}
	Diseñe un problema original de regresión lineal simple (distinto de calificaciones, paquetes, energía, aleaciones, automóviles o publicidad ya vistos): una aplicación de reparto a domicilio modela el tiempo de entrega ($Y$, en minutos) en función de la distancia recorrida ($X$, en km). Para una muestra de $n=12$ pedidos se resumen: $\bar{X}=5.0$ km, $\bar{Y}=22.0$ min, $S_{XY}=54.0$, $S_{XX}=30.0$, $\text{SCT}=150.0$. Ajuste el modelo por MCO, prediga el tiempo de entrega para un pedido de $8$ km, y calcule e interprete $R^2$.
\end{problema}

\subsection*{Soluciones detalladas}

% Recordar
\begin{solproblema}[prob:ab5bf4a]
	\begin{enumerate}
		\item $\bar{X}=66/10=6.6$, $\bar{Y}=675/10=67.5$. Con $\sum X_iY_i=4895.0$ y $\sum X_i^2=508.0$:
		\begin{align}
			S_{XY} &= 4895.0-10(6.6)(67.5) = 440.0, \qquad S_{XX} = 508.0-10(6.6)^2 = 72.4.
		\end{align}
		\begin{align}
			\hat{\beta}_1 = \frac{440.0}{72.4} \approx 6.0773, \qquad \hat{\beta}_0 = 67.5-(6.0773)(6.6) \approx 27.3898.
		\end{align}
		\item $\hat{Y} = 27.39+6.08X$.
		\item Para $X=7.5$: $\hat{Y} = 27.3898+6.0773(7.5) \approx 72.97$ puntos.
		\item $\text{SCR}=\hat{\beta}_1S_{XY}=(6.0773)(440.0)\approx2674.03$. Con $\sum Y_i^2=46825.0$: $\text{SCT}=46825.0-10(67.5)^2=1262.5$. Entonces $R^2=2674.03/1262.5\approx0.9972$ (99.72\%): el 99.72\% de la variabilidad en las calificaciones es explicada por las horas de estudio, una relación lineal casi funcional.
	\end{enumerate}
\end{solproblema}

% Comprender
\begin{solproblema}[prob:e5cf043]
	\begin{enumerate}
		\item $\hat{\beta}_0=5.0$: el costo base fijo de procesamiento, independiente del peso ($X=0$), es de \$5.00. $\hat{\beta}_1=2.5$: cada kilogramo adicional incrementa el costo de envío en \$2.50.
		\item Para $X=15.0$: $\hat{Y}=5.0+2.5(15.0)=42.50$ dólares.
		\item $\text{SCR}=R^2\times\text{SCT}=(0.85)(200.0)=170.0$. $\text{SCE}=\text{SCT}-\text{SCR}=200.0-170.0=30.0$.
		\item Con $\nu_E=20-2=18$: $\hat{\sigma}^2=30.0/18\approx1.6667$, $S_e=\sqrt{1.6667}\approx1.2910$ dólares.
	\end{enumerate}
\end{solproblema}

% Aplicar
\begin{solproblema}[prob:5a642fc]
	\begin{enumerate}
		\item $\hat{\beta}_1=2000.0/500.0=4.0$ MWh/°C. $\hat{\beta}_0=450.0-4.0(25.0)=350.0$. Recta: $\hat{Y}=350.0+4.0X$.
		\item $\text{SCR}=\hat{\beta}_1S_{XY}=(4.0)(2000.0)=8000.0$. $\text{SCE}=10000.0-8000.0=2000.0$. $R^2=0.800$.
		\item Con $\nu_E=13$: $\text{CME}=2000.0/13\approx153.85$, $SE(\hat{\beta}_1)=\sqrt{153.85/500.0}\approx0.5547$. El estadístico es $t=4.0/0.5547\approx7.211$. Como $|t|=7.211\gg t_{0.025,13}=2.160$, \textbf{se rechaza $H_0$}: existe una dependencia fuertemente significativa entre temperatura y consumo eléctrico.
	\end{enumerate}
\end{solproblema}

% Analizar
\begin{solproblema}[prob:e97453b]
	\begin{enumerate}
		\item Derivando $Q(\beta_0,\beta_1)=\sum(Y_i-\beta_0-\beta_1X_i)^2$ e igualando a cero:
		\begin{align}
			\frac{\partial Q}{\partial\beta_0} = -2\sum(Y_i-\beta_0-\beta_1X_i)=0, \qquad \frac{\partial Q}{\partial\beta_1} = -2\sum X_i(Y_i-\beta_0-\beta_1X_i)=0,
		\end{align}
		que al distribuir dan las \textbf{ecuaciones normales de Gauss}: $n\hat{\beta}_0+\hat{\beta}_1\sum X_i=\sum Y_i$ y $\hat{\beta}_0\sum X_i+\hat{\beta}_1\sum X_i^2=\sum X_iY_i$.
		\item De la primera ecuación: $\hat{\beta}_0=\bar{Y}-\hat{\beta}_1\bar{X}$. Sustituyendo en la segunda y despejando:
		\begin{align}
			\hat{\beta}_1\left[\sum X_i^2-n\bar{X}^2\right] = \sum X_iY_i-n\bar{X}\bar{Y} \implies \hat{\beta}_1 = \frac{\sum(X_i-\bar{X})(Y_i-\bar{Y})}{\sum(X_i-\bar{X})^2}.
		\end{align}
		\item Sustituyendo $Y_i=\beta_0+\beta_1X_i+\varepsilon_i$ en $\hat{\beta}_1$ y usando $\sum(X_i-\bar{X})=0$, $\sum(X_i-\bar{X})X_i=S_{XX}$:
		\begin{align}
			\hat{\beta}_1 = \beta_1 + \sum_{i=1}^n\left(\frac{X_i-\bar{X}}{S_{XX}}\right)\varepsilon_i.
		\end{align}
		Tomando esperanza con $E[\varepsilon_i]=0$: $E[\hat{\beta}_1]=\beta_1$. Para el intercepto, $E[\hat{\beta}_0]=E[\bar{Y}]-\bar{X}E[\hat{\beta}_1]=(\beta_0+\beta_1\bar{X})-\bar{X}\beta_1=\beta_0$. Ambos estimadores son insesgados.
	\end{enumerate}
\end{solproblema}

% Evaluar
\begin{solproblema}[prob:5c55133]
	\begin{enumerate}
		\item En $X_0=6.0\%$: $\text{SE}(\hat{\mu}_{Y|X_0})=4.2\sqrt{1/18+(3.0)^2/40.0}=4.2\sqrt{0.28056}\approx2.2246$ megaciclos. $\text{SE}(Y_{\text{pred}|X_0})=4.2\sqrt{1+0.28056}=4.2\sqrt{1.28056}\approx4.7528$ megaciclos.
		\item Con $\hat{Y}_0=120.0+15.0(6.0)=210.0$ y $t_{0.025,16}=2.120$:
		\begin{align}
			IC_{95\%}(\mu_{Y|X_0=6}) = 210.0\pm2.120(2.2246) = [205.28,\;214.72], \\
			IP_{95\%}(Y_{\text{pred}|X_0=6}) = 210.0\pm2.120(4.7528) = [199.92,\;220.08].
		\end{align}
		El intervalo de predicción ($\pm10.08$) es más del doble de ancho que el de confianza de la media ($\pm4.72$). Esto ocurre porque el IC estima solo dónde se sitúa el promedio de infinitos especímenes en las mismas condiciones —y se colapsa a ancho cero cuando $n\to\infty$—, mientras que el IP debe cubrir la dispersión de \emph{una observación individual futura}, que siempre incluye la varianza propia $\sigma^2$ de esa pieza. Por eso la banda de predicción nunca puede angostarse más allá de $\pm t_{\alpha/2}\sigma$, sin importar cuán grande sea la muestra.
	\end{enumerate}
\end{solproblema}

% Crear
\begin{solproblema}[prob:58d3723]
	\begin{align}
		\hat{\beta}_1 = \frac{S_{XY}}{S_{XX}} = \frac{54.0}{30.0} = 1.8 \text{ min/km}, \qquad \hat{\beta}_0 = 22.0-1.8(5.0) = 13.0 \text{ min}.
	\end{align}
	El modelo ajustado es $\hat{Y}=13.0+1.8X$. Para un pedido de $X=8$ km:
	\begin{align}
		\hat{Y} = 13.0+1.8(8) = 13.0+14.4 = 27.4 \text{ minutos.}
	\end{align}
	La suma de cuadrados de la regresión es $\text{SCR}=\hat{\beta}_1S_{XY}=(1.8)(54.0)=97.2$, así que:
	\begin{align}
		R^2 = \frac{\text{SCR}}{\text{SCT}} = \frac{97.2}{150.0} = 0.648 \text{ (64.8\%).}
	\end{align}
	El $64.8\%$ de la variabilidad en el tiempo de entrega se explica por la distancia recorrida; el $35.2\%$ restante se debe a otros factores no incluidos en el modelo (tráfico, hora del día, disponibilidad de repartidores, etc.).
\end{solproblema}
